\cvsubtitle{Undergraduate Projects}

\cventry{Characterizing Gravitational Wave Signals from Core-Collapse Supernovae}{Oct 2018 - Sept 2023}
{\vspace*{-3.5mm} \begin{itemize}
        \item{Modified and extended 1D core-collapse models written in FORTRAN.}
        \item{Modified and extended code to calculate gravitational wave eigenfrequencies given hydrodynamic data from a simulated supernova.}%%clown, glad u found this%%
        \item{Compared spacetime geometry assumed in eigenfrequency calculation code to that of our models, verifying their approximate equivalence numerically and with pen-and-paper.}
        \item{Applied active subspace analysis, a novel method of sensitivity analysis.}
        \item{"Poster-Blitz Advertisement" \& poster presentation at Fifty-One Ergs (FOE) in May 2019}
        \item{Created a system to automatically run a suite of models on multiple computers.}
        \item{Writing the final manuscript, with an anticipated publication date in Fall 2022.}
    \end{itemize}
}{with Professor Carla Fr\"ohlich, NC State}

\cventry{Parameter Inference with Core-Collapse Supernovae Models}{June 2020 - Sept 2023}{\vspace*{-3.5mm} \begin{itemize}
    \item Interdisciplinary astrophysics \& data science collaboration using machine learning and modern statistical methods to constrain key physical parameters like the nuclear equation of state.
    \item Identified and collected astronomical data that can be compared to the results of 1D core-collapse models.
    \item Performed simulations of core-collapse and post-processed them to extract relevant quantities for analysis in Meskhi, Wolfe et al. (2022), accepted to the Astrophysical Journal Letters.
\end{itemize}}{with collaborators at NC State, LANL, and University of Houston}

\cventry{Uncertainty Quantification for Core-Collapse Supernova Models}{Dec 2020 - May 2022}{\vspace*{-3.5mm} \begin{itemize}
    \item Independently ideated and proposed to address gaps in understanding of core-collapse simulation uncertainties, to better connect simulation results to astronomical observations.
    \item Built first machine learning autoencoder models to predict core-collapse supernova explosion energies with the Python Keras library, with little prior machine learning experience.
    \item Performed Bayesian parameter estimation using my machine learning model and Monte Carlo Markov Chain sampling to quantify uncertainty in explosion energy predictions.
\end{itemize}}{with Professors Carla Fr\"ohlich and Ralph Smith, NC State}

\cventry{Testing General Relativity with Gravitational Wave Signals using Hybrid Sampling}{June 2021 - May 2023}{\vspace*{-3.5mm} \begin{itemize}
    \item Implemented the TaylorF2 waveform approximant to modify post-Newtonian inspiral parameters.
    \item Implemented, from scratch, a unique method to combine multiple parameter estimation techniques in a new hybrid fashion.
    \item Integrated my sampling code with LIGO \code{BILBY} software and ran it on the LIGO Data Grid.
    %\item Discovered new limits on the detectability of deviations from general relativity in gravitational wave signals.
    \item Wrote two interim written reports and a final report for the Lab.
    \item Gave two interim progress talks and a final talk to the Lab.
\end{itemize}
}{with Dr. Colm Talbot, Caltech LIGO Astrophysics Lab \href{https://dcc.ligo.org/T2100339}{dcc.ligo.org/T2100339}}

\cventry{Identifying Black Hole Formation in Core-Collapse Supernova Simulations}{Sept 2021 - Dec 2021}{\vspace*{-3.5mm} \begin{itemize}
    \item Converted the metric of the general relativistic hydrodynamics code \code{AGILE} into the ADM 3+1 formalism, identifying the lapse $\alpha$ and shift $\beta^{i}$
    \item Rederived an expression for the expansion $H$, that characterizes apparent horizon formation, in the coordinate system of \code{AGILE}
    \item Calculated the expansion $H$ in Python to identify black hole formation in previously completed spherically-symmetric core-collapse supernova simulations
\end{itemize}
}{with Dr. Jonah Miller, LANL and Professor Carla Fr\"ohlich, NC State}

\cventry{Studying the Equation of State Dependence of Core-Collapse Supernovae Outcomes}{Aug 2019 - Aug 2021}{\vspace*{-3.5mm} \begin{itemize}
    \item Modified post-processing software used to calculate model explosion energy software
    \item Created suite of core-collapse models across six equations of state, supplying the simulation data used Ghosh et al. (2022).
\end{itemize}
}{with Somdutta Ghosh and Professor Carla Fr\"ohlich, NC State}

\cventry{Investigating Neutral Hydrogen in the Cold Accretion Disk around Sagittarius A*}{June 2019 - Oct 2020}{\vspace*{-3.5mm} \begin{itemize}
    \item Studied the description and observational consequences of the $\mathrm{H}30\alpha$ transition of hydrogen from $n = 31$ to $n = 30$, in the environment around a black hole, in the presence of significant \textit{non-blackbody} background radiation.
    \item Applied statistical mechanics, in a "pen-and-paper" manner, to understand how transition probabilities are modified by the non-blackbody background.
    \item Independently began this collaboration with Dr. Murchikova, during a conversation following her presentation at NC State.
\end{itemize}}{with Dr. Elena Murchikova, Member, Institute for Advanced Study}

\cventry{Applying Sensitivity Analysis Methods to Core-Collapse Supernova Lightcurves}{June 2020 - September 2020}{\vspace*{-3.5mm} \begin{itemize}
    \item Applied a novel method for sensitivity analysis, known as active subspace analysis to analyze the light curves resulting from 1D core-collapse supernovae models.
    \item Identified the ejected $^{56}\textrm{Ni}$ mass as the primary determinant of lightcurve properties.
    \item Contributed a subsection and appendix as second author on Curtis, Wolfe et al. (2021), in the Astrophysical Journal.
\end{itemize}}{with Dr. Sanjana Curtis, University of Amsterdam}

%\cventry{Optical Simulation of Organic Photovoltaic Devices}{Aug 2018 - July 2019}{\vspace*{-3.5mm} \begin{itemize}
%    \item Created optical simulations in Python to model light transmission through organic photovoltaic devices.
%    \item Attempted to calculate the change in exciton binding energy when exposed to an external electric field (via the Stark Effect).
%\end{itemize}
%}{with Research Professor Abay Dinku, NC State}